% Part: first-order-logic
% Chapter: natural-deduction
% Section: rules-and-proofs

\documentclass[../../../include/open-logic-section]{subfiles}

\begin{document}

\iftag{FOL}
      {\olfileid{fol}{ntd}{rul}}
      {\olfileid{pl}{ntd}{rul}}

\olsection{నియమాలు మరియు \usetoken{P}{derivation}}

\begin{explain}
సహజ నిగమన వ్యవస్థలు గణిత నిరూపణలో ఉపయోగించే అనౌపచారిక తార్కిక
ఆలోచనను దగ్గరగా అనుసరించడానికి ఉద్దేశించినవి (అందుకే అవి కొంతవరకు
“సహజమైనవి”). సహజ నిగమన నిరూపణలు పరికల్పనలతో మొదలవుతాయి. తరువాత
నిగమన నియమాలను వర్తింపజేస్తాం. \Intro{\lnot}, \Intro{\lif},
\iftag{FOL}{\Elim{\lor}, \Elim{\lexists}}{\Elim{\lor}} నిగమన
నియమాలు పరికల్పనలను “\tetoken{ఉపసంహరించిన}{!!{discharged}}” చేస్తాయి; స్పష్టత కోసం
\tetoken{ఉపసంహరించిన}{!!{discharged}} అయిన పరికల్పన గుర్తును ఆ నిగమనం పక్కన పెడతాం.
\end{explain}

\begin{defn}[పరికల్పన]
ఏ శాఖలోనైనా అన్నిటికంటే పైనున్న స్థానంలోని ఏ \tetoken{వాక్యమైనా}{!!{sentence}}
\emph{పరికల్పన} అంటాం.
\end{defn}

సహజ నిగమనంలోని \tetoken{వ్యుత్పత్తులు}{!!^{derivation}s} కొన్ని \tetoken{వాక్యాలు}{!!{sentence}s} వృక్షాలు.
అందులో అన్నిటికంటే పైనున్న \tetoken{వాక్యాలు}{!!{sentence}s} పరికల్పనలు; ఒక \tetoken{వాక్యం}{!!a{sentence}}
ఒకటి, రెండు లేదా మూడు ఇతర వాక్యాల కింద ఉంటే, అది నిగమన నియమం
ద్వారా సరిగ్గా రావాలి. నిగమనం పైనున్న \tetoken{వాక్యాలను}{!!{sentence}s} దాని
\emph{పూర్వాధారాలు} అనీ, కిందున్న \tetoken{వాక్యాన్ని}{!!{sentence}} దాని
\emph{నిష్కర్ష} అనీ అంటాం. ప్రతి \tetoken{తార్కిక సంయోజకానికి}{!!{operator}} ఒక ప్రవేశ నియమం,
ఒక తొలగింపు నియమం చొప్పున నియమాలు జతలుగా వస్తాయి. అవి నిష్కర్షలో
ఒక \tetoken{తార్కిక సంయోజకాన్ని}{!!a{operator}} ప్రవేశపెడతాయి లేదా నియమపు పూర్వాధారం నుంచి ఒక
\tetoken{తార్కిక సంయోజకాన్ని}{!!a{operator}} తొలగిస్తాయి. కొన్ని నియమాలు ఒక నిర్దిష్ట రకపు
పరికల్పనను \emph{\tetoken{ఉపసంహరించిన}{!!{discharged}}} చేయనిస్తాయి. ఏ నిగమనం ఏ పరికల్పనను
\tetoken{ఉపసంహరించిన}{!!{discharged}} చేస్తుందో సూచించడానికి పరికల్పనకూ నిగమనానికీ గుర్తులు
పెడతాం. పరికల్పనను ``$\Discharge{!A}{n}$'' అని రాసి దీన్ని సూచిస్తాం.
\sourcecorrection{OLTEND-001}{మూల నిర్వచనం ఇక్కడ “ఇతర సీక్వెంట్ల”ని
అంటుంది; ఈ వృక్షపు నోడ్‌లు వాక్యాలేనని అదే పేరా, పూర్వాధారాల నిర్వచనం
నిర్ణయిస్తున్నందున “ఇతర వాక్యాల”గా సరిచేశాం.}

% Only include this sentence is on of \land, lor, lif or lnot is defined.

\iftag{notprvNot,notprvAnd,notprvOr,notprvIf}{}%
	{$\land$, $\lor$, $\lif$, $\lnot$, $\lfalse$లలో కొన్ని నిర్వచితమైనవైనా, అన్ని \tetoken{తార్కిక సంయోజకాలకు}{!!{operator}s} నియమాలను పరిగణించడం ఆచారం.}



\end{document}
